\subsection{Module ranks and dimension-sixty Morita transport}
\begin{proposition}\label{prop:m3-a5-module-rank}
The ring $\sha{R8_3_307}$ admits no complex categorification.
\end{proposition}
\begin{proof}
Its positive dimensions, in source order, are
$(1,\sqrt {10},\sqrt {10},2\sqrt {10},5,3,3,4)$.
The integral basis $\{0,4,5,6,7\}$ is the neutral part of a faithful
$C_2$-grading and has dimension $60$. It has no nontrivial proper fusion
subring. Any neutral category $\mathcal D$ is therefore simple and integral,
so \cite[Theorem~9.12]{ENOSolvable} gives
$\mathcal D\simeq\operatorname{Rep}(A_5)$.

The odd component is an invertible $\mathcal D$-bimodule of rank three.
By the subgroup/projective-representation classification of
$\operatorname{Rep}(A_5)$-module categories
\cite[Theorem~2 and Remark~12]{OstrikModules}, its objects have projective
degrees $a(1,1,2)$: the exact action matrices have a one-dimensional
simultaneous dimension eigenspace with primitive vector $(1,1,2)$.
The certificate checks both its eigenvector equations and rank two of the
combined coefficient matrix. For the defining subgroup $H\leq A_5$,
$|H|=6a^2\mid60$, whence $a=1$ and $H\cong S_3$. A one-dimensional
projective representation trivializes the multiplier.

The dual category has one family of simples for each $H$-double coset,
parametrized by irreducible representations of its stabilizer.
The three double cosets have sizes $6,18,36$ and stabilizers $S_3,C_2,1$.
Its rank is consequently $3+2+1=6$. Invertibility requires its dual to be
$\mathcal D$, of rank five, a contradiction. The finite certificate checks
every order-six embedding through Sylow-three normalizers and the complete
double-coset calculation.
\end{proof}

\begin{proposition}[Pointed transport from dimension sixty]
\label{prop:m3-sixty-transport}
Let $\mathcal C$ be faithfully graded by a finite group, with a nonsolvable
neutral category of dimension $60$. Then $\mathcal C$ is group-theoretical
and integral. In particular $\sha{R8_3_2132}$ has no complex categorification.
\end{proposition}
\begin{proof}
By \cite[Theorem~9.16]{ENOSolvable}, the neutral category is weakly
group-theoretical. Choose a nilpotent Morita representative. The orders of
its extension groups multiply to $60$. If all these groups were solvable,
Morita invariance and \cite[Proposition~4.5]{ENOSolvable} would imply
solvability of the neutral category. Every group whose order is a proper
divisor of $60$ is solvable, while the only nonsolvable group of order $60$
is $A_5$. The representative must therefore be
$\operatorname{Vec}_{A_5}^{\omega}$, with arbitrary admissible associator
$\omega$.

Transport the entire grading under this neutral Morita equivalence, using
\cite[Lemma~3.4]{ENOSolvable}. Each component becomes an invertible
bimodule over $\operatorname{Vec}_{A_5}^{\omega}$, whose one-sided dual is
pointed. The pointed-dual criterion \cite[Theorem~3.4]{NaiduMorita}
requires the subgroup defining the module category to be abelian and
normal in $A_5$, hence trivial. Each component therefore has $60$ simples
and dimension $60$, so all its simples are invertible. The transported
category is pointed; the original category is group-theoretical and thus
integral. No assumption $\omega=1$ has been made.

For $\sha{R8_3_2132}$ the dimensions are
$(1,1,\sqrt {30},5,5,\sqrt {30},2,2)$.
Its integral neutral category $\mathcal D$ has dimension $60$, equals its
adjoint subcategory, and has two invertibles. If it were solvable, its last
cyclic-prime construction step could not be an extension, since
$\mathcal D_{\mathrm{ad}}=\mathcal D$. It would be a $C_2$-equivariantization,
by \cite[Proposition~4.4]{ENOSolvable}. Multiplication by the nontrivial
character has neutral orbits $\{0,1\},\{3,4\},\{6\},\{7\}$.
A fixed underlying simple yields two equivariant simples interchanged by
that character; a free underlying orbit yields one character-fixed simple
of twice its dimension. De-equivariantization would have dimensions
$1,1,1,1,1,5$, the near-group ring $\operatorname{NG}(C_5,4)$.
Its whole tensor matches the already excluded $\sha{R6_4_19}$; the
dimension-thirty Morita certificate permits only simple dimensions one
and two. Thus $\mathcal D$ is nonsolvable. The transport result makes the
whole category integral, contradicting its odd dimension $\sqrt {30}$.
\end{proof}
The source-bound replay checks all dimension and grading equations,
adjoint closure, character orbits, the complete lower-rank tensor match
and the finite normal-subgroup premise. This is a new application and
unification of the dimension-sixty dichotomy and integrality reductions;
it does not identify an arbitrary categorification from its ring alone.
